\documentclass[11pt,a4paper]{article}

% ============================================================
% Packages
% ============================================================
\usepackage[utf8]{inputenc}
\usepackage[T1]{fontenc}
\usepackage{lmodern}
\usepackage{microtype}
\usepackage[margin=2.5cm]{geometry}

\usepackage{amsmath,amssymb,amsthm}
\usepackage{mathtools}
\usepackage{xspace}
\usepackage{xcolor}
\usepackage{booktabs}
\usepackage{longtable}
\usepackage{float}
\usepackage{array}
\usepackage{enumitem}
\usepackage{listings}
\usepackage{tikz}
\usetikzlibrary{positioning,arrows.meta}

\usepackage[colorlinks=true,linkcolor=blue!55!black,citecolor=green!45!black,urlcolor=blue!55!black]{hyperref}
\usepackage[capitalize,nameinlink]{cleveref}

\newcolumntype{L}[1]{>{\raggedright\arraybackslash}p{#1}}

% ============================================================
% Theorem environments
% ============================================================
\theoremstyle{plain}
\newtheorem{theorem}{Theorem}[section]
\newtheorem{claim}[theorem]{Claim}
\theoremstyle{definition}
\newtheorem{definition}[theorem]{Definition}
\newtheorem{remark}[theorem]{Remark}

% ============================================================
% Rocq listing style
% ============================================================
\definecolor{coqkw}{rgb}{0.30,0.10,0.55}
\definecolor{coqcom}{rgb}{0.30,0.45,0.30}
\definecolor{coqstr}{rgb}{0.55,0.20,0.10}
\lstdefinelanguage{Coq}{
  morekeywords={Axiom,Definition,Lemma,Theorem,Corollary,Record,Inductive,
    Fixpoint,Variant,Context,Variable,Parameter,Notation,Proof,Qed,Defined,
    forall,fun,match,with,end,let,in,if,then,else,Module,Section,End,Import,
    Require,From,bool,nat,package,game,perfect,return,as,of,Type,Prop},
  sensitive=true,
  morecomment=[s]{(*}{*)},
  morestring=[b]",
}
\lstset{
  language=Coq,
  basicstyle=\ttfamily\footnotesize,
  keywordstyle=\color{coqkw}\bfseries,
  commentstyle=\color{coqcom}\itshape,
  stringstyle=\color{coqstr},
  columns=fullflexible,
  keepspaces=true,
  showstringspaces=false,
  breaklines=true,
  breakatwhitespace=true,
  xleftmargin=1em,
  frame=single,
  rulecolor=\color{black!15},
  numbers=left,
  numberstyle=\tiny\color{black!45},
  numbersep=6pt,
  literate=
    {≈}{{$\approx$}}1 {∘}{{$\circ$}}1 {≤}{{$\leq$}}1 {≥}{{$\geq$}}1
    {→}{{$\rightarrow$}}1 {←}{{$\leftarrow$}}1
    {⦃}{{$\{\!|$}}1 {⦄}{{$|\!\}$}}1 {∑}{{$\sum$}}1
    {×}{{$\times$}}1 {∈}{{$\in$}}1 {⊢}{{$\vdash$}}1
    {∐}{{$\sqcup$}}1 {α}{{$\alpha$}}1 {λ}{{$\lambda$}}1
    {⋊}{{$\rtimes$}}1 {∗}{{$*$}}1 {∀}{{$\forall$}}1,
}
\DeclareRobustCommand{\coq}[1]{\texttt{\detokenize{#1}}}

% ============================================================
% Notation
% ============================================================
\newcommand{\SSProve}{\textsf{SSProve}\xspace}
\newcommand{\Rocq}{\textsf{Rocq}\xspace}
\newcommand{\WOTSC}{\textsf{WOTS+C}\xspace}
\newcommand{\Adv}{\mathsf{Adv}}
\newcommand{\AdvOf}{\mathsf{AdvOf}}
\newcommand{\A}{\mathcal{A}}
\newcommand{\B}{\mathcal{B}}
\newcommand{\PKseed}{\ensuremath{\mathsf{PKseed}}\xspace}
\newcommand{\SKseed}{\ensuremath{\mathsf{SKseed}}\xspace}
\newcommand{\prf}{\textsc{prf}\xspace}
\newcommand{\smud}{\textsc{sm-ud}\xspace}
\newcommand{\smtcr}{\textsc{sm-tcr}\xspace}
\newcommand{\smpre}{\textsc{sm-pre}\xspace}
\newcommand{\stcrp}{\textsc{s-tcr-p}\xspace}
\newcommand{\epsprf}{\varepsilon_{\mathsf{prf}}}
\newcommand{\epsud}{\varepsilon_{\mathsf{ud}}}
\newcommand{\epsgr}{\varepsilon_{\mathsf{grind}}}
\newcommand{\epstpk}{\varepsilon_{\mathsf{tcr,pk}}}
\newcommand{\epstch}{\varepsilon_{\mathsf{tcr,chain}}}
\newcommand{\epspre}{\varepsilon_{\mathsf{pre}}}
\newcommand{\pathcode}[1]{\path{#1}}

\title{\textbf{Reviewing the Machine-Checked Security Proof of \WOTSC}\\[4pt]
       {\large A Guide for Cryptographers}}
\author{}
\date{7 August 2026\\[2pt]
      \small Reviewed commit \texttt{be8d966660973f1367f8d3cc66886472649e59e3}}

\begin{document}
\maketitle

\begin{abstract}
This guide explains how to review the machine-checked \WOTSC security proof formalized in \Rocq using \SSProve, a Rocq library for game-based cryptographic proofs.
For each adversary in the formal \WOTSC experiment, the theorem constructs adversaries against six assumption games and bounds the original advantage by the sum of their advantages.
The proved experiment fixes the signed digest before the adversary sees the public key.
This known-message notion is sufficient when the security proof of a larger hash-based signature scheme fixes the internal digest before invoking the \WOTSC result.
Rocq checks the reduction, while a final corollary substitutes assumed upper bounds for the six game advantages.
The guide explains how to inspect the experiment, reduction, assumption games, and unproved dependencies.
It also compares the formal games with published QROM bounds: using those bounds requires game-correspondence and resource arguments outside the checked proof, and the tight bound relevant to the length-preserving preimage game is conjectural.
\end{abstract}

\tableofcontents
\newpage

\section{What the theorem proves}
\label{sec:claim}

The formal proof establishes a reduction from forging a \WOTSC signature to breaking one of six precisely defined security games for the functions used by \WOTSC.
This section states that result without relying on Rocq or \SSProve syntax.

\subsection{The \WOTSC signature scheme}

\WOTSC signs a digest by revealing intermediate values from a set of one-way hash chains.
Each chain begins at a secret value and ends after a fixed number of calls to the chain hash.
Key generation derives the secret chain starts from \(\SKseed\), advances every chain to its endpoint, and compresses all endpoints into one value.
The public key consists of that compressed value and the public seed \(\PKseed\) used to domain-separate the hash calls.

Before choosing positions in the chains, the signer combines the digest with a counter and hashes the pair.
The hash output determines one position in each chain.
The signer tries counters in increasing order until the sum of these positions equals the value required by \WOTSC, or stops when it reaches the counter limit \(\mathsf{cap}\).
A signature contains the value at each selected chain position and the counter that selected those positions.
The Rocq files call this bounded counter search \emph{grinding}.

The verifier derives the same positions from the digest and counter.
It advances every signature value from its selected position to the end of its chain, compresses the resulting endpoints, and compares the result with the public key.
The security proof concerns these key-generation, signing, and verification algorithms.

\subsection{The attack experiment}

The theorem analyzes the following interaction.

\begin{enumerate}[leftmargin=1.6em,itemsep=2pt]
\item The adversary chooses a digest \(m\) to be signed.
\item The challenger generates a fresh \WOTSC key pair, signs \(m\), and gives the public key and signature to the adversary.
\item The adversary submits candidate pairs \((m^\star,\sigma^\star)\).
\item A candidate is a forgery if \(m^\star\neq m\) and verification accepts \(\sigma^\star\) for \(m^\star\).
\end{enumerate}

The digest \(m\) is therefore fixed before the adversary sees the public key.
The Rocq files call this a one-time known-message experiment.
It is not standard one-time EUF-CMA, in which the adversary sees the public key before choosing its signing query.
This known-message notion is sufficient when \WOTSC is used inside a larger hash-based signature scheme whose security proof fixes the internal digest before invoking the \WOTSC theorem.
The exact treatment of repeated candidate submissions and signing failure is part of the formal game reviewed in \cref{sec:notion}.

\subsection{The proved reduction}

For each adversary in the formal experiment, the proof constructs adversaries against six games: one for secret-key derivation, one for the distribution of internal chain values, one for the digest-and-counter hash, two target-collision games, and one preimage game.
In the following display, \(\mathsf{KOTS\_vc}\) names the attack experiment just described, and \(\AdvOf(\mathcal{G},\A)\) denotes the advantage assigned to adversary \(\A\) by game \(\mathcal{G}\).
\Cref{sec:ssprove-adv} gives the exact \SSProve definition of this advantage.
Writing the six assumption bounds as \(\varepsilon\)-terms, the theorem gives the following inequality:
\begin{equation}
\label{eq:security-bound}
  \AdvOf(\mathsf{KOTS\_vc},\A)
  \leq 2\epsprf+2(w-2)\epsud+\epsgr+\epstpk+\epstch+\epspre.
\end{equation}
Here \(\epsprf\) bounds the key-derivation game, \(\epsud\) bounds the chain-value undetectability game, and \(\epsgr\) bounds the digest-and-counter collision game.
The remaining terms bound target collisions in public-key compression, target collisions in the chain hash, and preimages of chain-hash outputs.

The theorem is pointwise: for a particular adversary \(\A\), each term on the right is the advantage of a specific adversary constructed from \(\A\).
\Cref{sec:reduction} shows the exact Rocq statement and the constructed adversaries.

\section{Scope and limitations}
\label{sec:limits}

The reduction in \cref{sec:claim} is machine checked, but several broader security claims require arguments that are not present in the Rocq theorem.
This section collects those limitations in one place.
Later sections refer back to this section when a line of code determines one of them.

\begin{table}[H]
\centering
\caption{Claims that require an argument beyond the checked reduction.}
\label{tab:limits}
\small
\begin{tabular}{@{}L{0.29\textwidth}L{0.63\textwidth}@{}}
\toprule
Broader claim & What is missing \\
\midrule
Time-bounded concrete security & \SSProve does not assign a cost to the adversary or the constructed reductions, and the six concrete axioms quantify over every package satisfying \coq{ValidPackage}. This predicate checks interfaces and memory use but does not impose a running-time bound. \\
Quantum or QROM security & The games and reductions have classical semantics, not quantum or QROM semantics. \\
Use of published QROM bounds & No checked theorem relates the formal games to the published experiments, so their bounds cannot be substituted directly. In addition, the tight published bound relevant to the length-preserving \smpre game depends on a conjectured multi-target decisional second-preimage bound. \\
\bottomrule
\end{tabular}
\end{table}

Published analyses can supply numerical bounds after a separate argument relates their experiments and resource bounds to the formal games reviewed here.
The attack experiment also aborts when the bounded counter search fails.
This is part of the proved security notion rather than a gap in the reduction, but a claim about an experiment that reveals signing failure would require a separate failure analysis.


\section{Trusted computing base, assumptions, and review obligations}
\label{sec:trust}

A reviewer must separate three questions that are easy to conflate.
First, did Rocq verify that the theorem follows from the formal definitions and listed assumptions?
Second, does the formal attack experiment express the intended security notion?
Third, do the assumption games match the properties for which external bounds are available?

The \emph{trusted computing base} (TCB) comprises the components that must behave correctly for Rocq's acceptance of a proof to be meaningful.
The Rocq kernel checks the completed proof, so this result relies on the kernel and on the software and hardware that execute it.
The Rocq kernel also checks the \SSProve and MathComp lemmas used by the proof.
Their definitions determine the meaning of packages, distributions, and advantages used in the theorem.

When Rocq checks a theorem, it also checks the proofs of the lemmas used by that theorem.
Following these uses transitively eventually reaches constants that were declared without proofs, such as an \coq{Axiom} or an imported admitted lemma.
This guide calls such a constant an \emph{unproved dependency} of the theorem.
The command \coq{Print Assumptions} lists these dependencies.

The unproved dependencies of the abstract reduction include proof irrelevance, choice, propositional and functional extensionality, axioms for real numbers, and an admitted double-series interchange lemma from an upstream real-analysis library.
A separate concrete corollary additionally depends on six cryptographic assumptions declared in \pathcode{proof/ssprove/axioms.v}.

The command \coq{make audit} rebuilds the proofs, rejects unfinished proofs such as \coq{Admitted}, checks that assumptions are declared only in the designated assumption files, and runs \coq{Print Assumptions} for the principal theorems.
It fails if these theorems acquire an unproved dependency outside the project's recorded set.
Passing the audit does not show that a cryptographic assumption is true, that the attack experiment expresses the intended security notion, or that a published bound applies to a formal game.
\Cref{sec:review-start} explains how to run this check.

\section{Check the proof with Rocq}
\label{sec:review-start}

All code line numbers in this guide refer to the commit printed on the title page.
From the repository root, enter the \SSProve development shell and run the proof target:

\begin{verbatim}
nix develop .#ssprove
make ssprove-prove
\end{verbatim}

The target succeeds only if the Rocq kernel accepts the proof files and their dependencies.
Then run the dependency audit:

\begin{verbatim}
make audit
\end{verbatim}

The final report combines the dependencies of several audited theorems.
It includes the six cryptographic assumptions because the concrete corollary uses them.
The individual \coq{Print Assumptions} queries are listed in \pathcode{proof/assumptions/ssprove.v}; the query for the abstract reduction should report only the classical dependencies described in \cref{sec:trust}.

\section{How to read the Rocq and SSProve files}
\label{sec:rocq-code}

This section introduces the Rocq and \SSProve notation needed to read the theorem and game definitions.
\Rocq is a proof assistant: its kernel checks that a theorem follows from the definitions, hypotheses, and unproved assumptions on which it depends.
\SSProve is a Rocq library for game-based cryptographic proofs.
It provides a language for probabilistic programs with oracles and private state, together with definitions of their semantics and distinguishing advantage.
In this proof, the attack experiment, assumption games, adversaries, and reductions are \SSProve programs, and Rocq checks the proofs about those programs.
The reviewer does not need to read the proof scripts or repeat the proof step by step.
The later sections follow the security statement from left to right: the attack experiment, the abstract reduction, the assumption games, and the unproved assumptions.

\subsection{The Rocq declarations used in this guide}

Every Rocq expression has a type.
A \coq{Definition} names an expression, and a \coq{Record} declares a type with named fields.
A \coq{Lemma} or \coq{Theorem} states a proposition followed by a proof that the kernel checks.
An \coq{Axiom} or \coq{Parameter} declares a constant without a definition or proof.
The command \coq{Print name} displays a declaration after Rocq has filled in omitted arguments, while \coq{Print Assumptions name} displays the unproved constants on which it depends.

A theorem has the following shape:

\begin{lstlisting}[numbers=none]
Theorem name (x : T) (H : P x) : Q x.
Proof.
  (* proof script *)
Qed.
\end{lstlisting}

The binder \coq{(x : T)} introduces a parameter \coq{x} of type \coq{T}, and \coq{(H : P x)} introduces a named hypothesis.
The declaration says that \(Q(x)\) holds for every \(x\) of type \(T\) satisfying hypothesis \(P(x)\).
The reviewer reads the declaration through the conclusion after the colon and may skip the commands between \coq{Proof} and \coq{Qed}.
The keyword \coq{Admitted} can close a theorem without a proof; such a theorem appears as an unproved dependency.

A \coq{Section} groups declarations, and each \coq{Context} variable used by a theorem becomes a universally quantified theorem parameter when the section closes.
For example, a theorem written under \coq{Context (P : wc_params)} is a theorem for every parameter record \coq{P}, not a theorem for one implicit global instance.
Rocq can fill in arguments and type conversions that do not appear in the file, so the output of \coq{Print} is authoritative when the written declaration is ambiguous.

Rocq distinguishes propositions from Boolean values.
The type \coq{Prop} is used for propositions in theorem statements, while \coq{bool} contains the executable values \coq{true} and \coq{false}.
Game code uses Boolean comparisons, while theorem statements use propositions such as real-number inequalities.

{\small
\begin{longtable}{@{}L{0.32\textwidth}L{0.60\textwidth}@{}}
\caption{The Rocq and MathComp notation used in later listings.}
\label{tab:rocq-notation}\\
\toprule
Notation & Meaning \\
\midrule
\endfirsthead
\toprule
Notation & Meaning \\
\midrule
\endhead
\bottomrule
\endfoot
\coq{fun x => e} & The anonymous function that maps \(x\) to expression \(e\). \\
\coq{let x := e in k} & Bind the deterministic value of \(e\) to \(x\) while evaluating \(k\). \\
\coq{match r with ... end} & Inspect the constructor of \(r\) and execute the corresponding branch. \\
\coq{option T} & The type of an optional \(T\) value. \\
\coq{Some x} & An optional value that contains \(x\). \\
\coq{None} & An optional value that contains no value. \\
\coq{X * Y} & Rocq's type for a pair whose components have types \(X\) and \(Y\). \\
\coq{(x, y)} & A pair containing \(x\) and \(y\). \\
\(\coq{X}\mathbin{\times}\coq{Y}\) & \SSProve's type for a pair whose components have types \(X\) and \(Y\). \\
\coq{'(x, y)} & A pattern that binds the two components of a pair to \(x\) and \(y\). \\
\coq{x.1} & The first component of pair \(x\). \\
\coq{x.2} & The second component of pair \(x\). \\
\(\coq{X}\mathbin{\sqcup}\coq{Y}\) & A disjoint-union type whose values contain either an \(X\) or a \(Y\). \\
\coq{inl x} & A disjoint-union value containing \(x\) from the left-hand type. \\
\coq{inr y} & A disjoint-union value containing \(y\) from the right-hand type. \\
\texttt{\{| f := e; ... |\}} & Construct a record by assigning a value to each named field. \\
\coq{[::]} & The empty sequence. \\
\coq{x :: xs} & The sequence obtained by prepending \(x\) to \(xs\). \\
\coq{rcons xs x} & The sequence obtained by appending \(x\) to \(xs\). \\
\coq{n.-2} & Natural-number subtraction truncated at zero. \\
\texttt{\textbackslash sum\_(0 <= s < n) e} & The finite sum of expression \(e\) over natural numbers \(s\) from zero through \(n-1\). \\
\coq{forall x, P} & Universal quantification: proposition \coq{P} holds for every value \coq{x}. \\
\coq{X -> Y} & The type of functions from \(X\) to \(Y\). \\
\coq{P -> Q} & When \(P\) and \(Q\) are propositions, the same notation is logical implication. \\
\texttt{(x < y)\%N} & Compare the natural numbers \(x\) and \(y\). \\
\texttt{(x == y)\%B} & Test \(x\) and \(y\) for equality and return a Boolean. \\
\end{longtable}
}

These constructs are enough to read the definitions and theorem statements in the rest of the guide.
Proof scripts contain many more tactics and library lemmas, but their meaning is not needed to audit the cryptographic statement.

\subsection{Probabilistic programs and oracle packages}
\label{sec:ssprove-semantics}

\SSProve represents a cryptographic game as a probabilistic program with oracles and private state~\cite{SSProve,SSProveJournal}.
It represents the adversary as another probabilistic program whose oracle calls are connected to the game.
The composition \(\A\circ G\) means that adversary \(\A\) runs with its imported oracles supplied by game \(G\).

An interface is a finite set of typed oracle identifiers.
The declaration \coq{[ O ] : { X ~> Y }} means that oracle \coq{O} accepts an input in \coq{X} and returns a value in \coq{Y}.
A \emph{package} implements its exported oracles, may call imported oracles, and owns typed private memory locations.
Both games and adversaries are packages.

A randomized computation returning a value of type \coq{T} has type \coq{dist T}.
The notations \coq{'bool}, \coq{'nat}, \coq{'list T}, and \coq{'option T} are the corresponding \SSProve datatypes.

\begin{table}[H]
\centering
\caption{The \SSProve program notation used in the game listings.}
\label{tab:notation}
\small
\begin{tabular}{@{}L{0.32\textwidth}L{0.60\textwidth}@{}}
\toprule
Notation & Operational meaning \\
\midrule
\texttt{x }\(\leftarrow\)\texttt{ d ;; k} & Run computation \(d\), bind its result to \(x\), and continue with \(k\). \\
\coq{ret x} & Return \(x\) without another randomized step. \\
\texttt{\{code c\}} & Use probabilistic computation \(c\) as the value of a record field. \\
\coq{mkloc n x} & Define a typed private-memory location with numeric name \(n\) and initial value \(x\). \\
\texttt{[package [fmap l1; ...]; ...]} & Define a package whose private state contains the listed memory locations. \\
\texttt{x }\(\leftarrow\)\texttt{ get l ;; k} & Read private location \(l\), bind its current value to \(x\), and continue with \(k\). \\
\texttt{\#put l := x ;; k} & Write \(x\) to private location \(l\) and continue with \(k\). \\
\texttt{\#assert b ;; k} & Continue with \(k\) if \(b=\coq{true}\), and abort otherwise. \\
\coq{fail} & Abort the current execution without returning a result. \\
\coq{getSome l} & Read optional location \(l\), return its contained value, or abort if it contains \coq{None}. \\
\coq{x == y} & Test \(x\) and \(y\) for equality and return a Boolean. \\
\coq{x != y} & Test \(x\) and \(y\) for inequality and return a Boolean. \\
\coq{~~ b} & Compute Boolean negation. \\
\coq{b1 && b2} & Compute Boolean conjunction. \\
\coq{setm m k v} & Update finite map \(m\) by associating key \(k\) with value \(v\). \\
\coq{getm m k} & Return the value associated with key \(k\) in finite map \(m\). \\
\bottomrule
\end{tabular}
\end{table}

The following excerpt illustrates how these forms compose in an oracle:

\begin{lstlisting}[numbers=none]
[ KFORGE ] : { Msg S × Sig S ~> 'bool } '(m, s) {
  pk ← getSome kots_pk_loc ;;
  log ← get kots_log_loc ;;
  if b then
    ret (verify S pk m s && (m \notin log))
  else
    ret false
}
\end{lstlisting}

The header declares an oracle named \coq{KFORGE} that accepts a digest-signature pair and returns a Boolean.
The pattern \coq{'(m, s)} names the two components of the input.
The first two commands read the stored public key and signed-digest log; \coq{getSome} aborts if no public key has been stored.

The aborting operations in the table require special care.
\SSProve uses subprobability distributions, so an aborted execution loses its probability mass instead of returning a failure value.
The remaining distribution is not normalized by conditioning on executions that continue.
This semantics explains the signing-failure limitation in \cref{sec:limits}.

\subsection{Real and ideal games}

The files express each security experiment as two games with the same interface.
In the forgery experiment, the real game checks the candidate and the ideal game rejects it; later sections state the corresponding difference for each assumption game.
Each pair is written as one definition indexed by a Boolean \coq{b}, where \coq{true} selects the real game and \coq{false} selects the ideal game.

This real-and-ideal presentation is \SSProve's way of expressing the advantage used by the reduction.
The later review must still check that the two games differ only where the intended security event is scored.
If an assertion or state lookup occurs in only one branch, an invalid oracle sequence may distinguish the games without producing a cryptographic forgery.

\subsection{Advantage and the \coq{ValidPackage} condition}
\label{sec:ssprove-adv}

\SSProve defines advantage as follows:

\begin{lstlisting}[numbers=none,mathescape=true,caption={The elaborated \SSProve advantage definition.},label={lst:adv}]
Definition Adv (G G' A : nom_package) : R :=
  locked `|distr.mu (Pr' (A $\circ$ G)) true
          - distr.mu (Pr' (A $\circ$ G')) true|.
Notation AdvOf GG A := (Adv (GG true) (GG false) A).
\end{lstlisting}

The expression \(\mathsf{Pr}'(\A\circ G)\) is the output distribution of adversary \(\A\) running with game \(G\), and \coq{distr.mu ... true} is the probability assigned to output \coq{true}.
Thus \(\Adv(G,G',\A)\) is the absolute difference between the probabilities that \(\A\) outputs \coq{true} in the two games.
The \coq{locked} wrapper controls when Rocq unfolds the definition during a proof but does not change its value.
The notation \(\AdvOf(\mathcal{G},\A)\) specializes this definition to a Boolean-indexed pair \(\mathcal{G}(\mathsf{true})\) and \(\mathcal{G}(\mathsf{false})\).

The Rocq type \coq{nom_package} is \SSProve's representation of a package whose oracles and memory locations have names.
The predicate \coq{ValidPackage} checks that such a package has the declared interfaces and uses its memory legally.
It does not restrict the order in which the adversary calls oracles, and it does not impose a running-time bound.

\section{Review the exact security notion}
\label{sec:notion}

The left side of the theorem is the advantage of an adversary against \coq{KOTS_vc_game} from \pathcode{proof/ssprove/lib/sig/kots.v}.
The interface fixes what the adversary may ask, and the two oracle bodies fix the order of events and the forgery condition.

\subsection{Interface and private state}

\Cref{lst:kots-interface} declares two oracles.

\lstinputlisting[
  caption={The two-oracle interface used by the \WOTSC theorem.},
  label={lst:kots-interface},
  firstline=35,lastline=45,firstnumber=35
]{../proof/ssprove/lib/sig/kots.v}

The \coq{COMMIT} input contains a scheme parameter and the digest to sign.
For the \WOTSC instance, the scheme parameter has type \coq{unit}, so the digest is the only input value chosen by the adversary.
The output contains the public key and an optional signature.
Because the public key is an output rather than an input, the adversary must supply the digest before learning the public key.

The \coq{KFORGE} oracle accepts a candidate digest and signature and returns a Boolean.
The package stores a public key and a sequence of signed digests in two private memory locations.
The expression \coq{mkloc 115 None} defines the public-key location, with its type fixed by the annotation \coq{option (Pk S)} and its initial value set to \coq{None}.
Similarly, \coq{mkloc 116 [::]} defines the signed-digest log with an initially empty sequence.
The numbers 115 and 116 are unique location names and have no cryptographic meaning.

\lstinputlisting[
  caption={The complete \coq{KOTS_vc_game} definition.},
  label={lst:kots-game},
  firstline=88,lastline=119,firstnumber=88
]{../proof/ssprove/lib/sig/kots.v}

\subsection{The \coq{COMMIT} oracle}

Lines 92 and 93 require the public-key location to be empty.
Once a \coq{COMMIT} call succeeds, another call aborts at this assertion.
Line 94 generates a fresh key pair, and line 95 signs the input digest \(m\).

If signing succeeds, lines 98 to 100 store the public key, record \(m\), and return \((pk,\coq{Some}\;\sigma)\).
If signing fails, line 101 aborts without returning a value or changing either location.
Thus the optional-signature component of the return type is always \coq{Some} in any execution that continues.
As explained in \cref{sec:limits}, the abort does not reveal a signing-failure result to the adversary.

\subsection{The \coq{KFORGE} oracle}

Lines 106 to 107 retrieve the stored public key and signed-digest log before examining the game bit \(b\).
A call to \coq{KFORGE} before a successful \coq{COMMIT} therefore aborts at the same point in both games.
This placement prevents an out-of-order call from distinguishing the games.

When \(b=\mathsf{true}\), lines 108 and 109 return
\[
  \mathsf{Verify}(pk,m^\star,\sigma^\star)\land(m^\star\notin[m]).
\]
The log contains only \(m\), so the second conjunct is equivalent to \(m^\star\neq m\).
This branch returns \coq{true} exactly when \(\sigma^\star\) verifies for a digest different from the signed digest.

When \(b=\mathsf{false}\), line 111 returns \coq{false} after the same state lookups.
The two games therefore differ only in the reply to a candidate submitted after a successful commitment.
If the adversary submits one candidate and outputs the oracle's reply, its advantage equals the probability that signing succeeds and that candidate is a valid forgery.
This probability is not conditioned on signing success: an execution that aborts during signing produces no final output.

The interface also permits several \coq{KFORGE} calls, so the adversary may test several candidates.
This does not provide access to a secret verification procedure because the verification algorithm and public key are already public.
The usual experiment in which the adversary returns one forgery is the special case with one \coq{KFORGE} call.

\section{Review the abstract reduction theorem}
\label{sec:reduction}

The main theorem in \pathcode{proof/ssprove/wots/theorem.v} bounds the attack experiment by advantages in the assumption games.
Each term on the right is the advantage of a specific package that the reduction constructs from the original adversary.
\prf stands for \emph{pseudorandom function}.
In \smud, \smtcr, and \smpre, the prefix \emph{SM} stands for \emph{single-function, multi-target}; \emph{UD}, \emph{TCR}, and \emph{PRE} stand for \emph{undetectability}, \emph{target-collision resistance}, and \emph{preimage resistance}.
In \stcrp, \emph{S} stands for \emph{special}, \emph{TCR} for \emph{target-collision resistance}, and the final \emph{P} for the output-predicate parameter.

\lstinputlisting[
  caption={The abstract \WOTSC security reduction.},
  label={lst:abstract-theorem},
  firstline=286,lastline=302,firstnumber=286
]{../proof/ssprove/wots/theorem.v}

\subsection{Read the theorem inputs and the left-hand side}

The names \(P\), \(K\), \coq{PRF}, and \coq{dK} are declared earlier in the file by four \coq{Context} commands.
As explained in \cref{sec:rocq-code}, Rocq adds these section parameters as universal quantifiers when the section closes.
They represent the \WOTSC parameter record, the secret-seed type, the key-derivation function, and the secret-seed sampler.
Write \(\ell=\coq{wcp_chain_count P}\) for the number of chains and \(w=\coq{wcp_w P}\) for the base used to encode a digest as chain positions.
The hypothesis \coq{Hrange} requires every digest-and-counter hash output to lie between zero and \(w^\ell-1\).
The exponent \(w^\ell\) is the number of length-\(\ell\) strings of base-\(w\) digits, so this condition makes the conversion from a hash output to such a string injective.
A separate Rocq lemma proves \coq{Hrange} for the parameters used by the final theorem.

Lines 288 and 289 quantify over an adversary package \(\A\).
The hypothesis \coq{VA} states that \(\A\) satisfies the structural \coq{ValidPackage} check described in \cref{sec:ssprove-adv}.
The expression \coq{loc A} denotes the package's private memory locations, and \coq{A_export} is the fixed interface through which it returns its final Boolean output.
The expression \coq{I_KOTS (shrincs_wots_scheme P)} is the \coq{COMMIT}/\coq{KFORGE} interface reviewed in \cref{sec:notion}.

Line 290 is the left-hand side of the inequality.
The record \coq{shrincs_wots_keyed P K PRF dK} supplies the key-generation, signing, and verification algorithms, with chain starts derived using \coq{PRF} and a secret seed sampled from \coq{dK}.
The constructor \coq{KOTS_vc_game} places those algorithms in the attack experiment from \cref{sec:notion}.
Thus line 290 is \(\A\)'s advantage in that experiment.

\subsection{Read the assumption terms}

Every term on lines 291 to 302 has the form \(\AdvOf(G,\A\circ R)\).
Here \(G\) is an assumption game and \(R\) is a reduction package defined in the proof.
The composed package \(\A\circ R\) attacks \(G\): \(R\) presents the \WOTSC oracles expected by \(\A\), while answering those calls through the oracles provided by \(G\).
Each \(R\) is an explicit Rocq definition, so the theorem identifies the reduction used for every assumption term.

In lines 292 and 293, \coq{s} is the summation index satisfying \(0\leq s<w-2\) and identifies one step between neighboring chain-value hybrids.
At step \coq{s}, the \smud game moves the hidden random value from position \coq{s} to position \coq{s + 1} in each chain whose corresponding base-\(w\) digest digit is at least \coq{s + 2}; all other chains remain unchanged.
The \coq{true} argument to \coq{R_ud} says that the simulated \coq{KFORGE} oracle checks candidate forgeries.
It does not select one side of the \smud game: \coq{AdvOf} compares both sides of that game.
The second \smud sum on lines 300 and 301 uses \coq{false} because its simulated \coq{KFORGE} oracle always rejects candidates.

The game constructors take records and parameter values that select the challenged function and configure the experiment.
\Cref{tab:theorem-names} gives the meaning of each such name in the order in which it first appears in the theorem.
The next section defines the corresponding experiments and their winning conditions.

{\small
\begin{longtable}{@{}L{0.32\textwidth}L{0.60\textwidth}@{}}
\caption{Names used to instantiate the assumption games in \cref{lst:abstract-theorem}.}
\label{tab:theorem-names}\\
\toprule
Name & Meaning in the theorem \\
\midrule
\endfirsthead
\toprule
Name & Meaning in the theorem \\
\midrule
\endhead
\bottomrule
\endfoot
\coq{khf_prf P K PRF dK} & The keyed function that derives one chain start from the secret seed, together with its input and output types and samplers. \\
\coq{wcp_chain_count P} & The number of hash chains, used here as the number of distinct \prf inputs or as the number of targets. \\
\coq{thf_chain P} & The chain hash, packaged with the types and samplers required by the tweakable-hash games. \\
\coq{coll_chain P} & Access to public-key compression and the digest-and-counter hash while the chain hash is being challenged. \\
\coq{wcp_w P} & The base used to convert a digest into chain positions. \\
\coq{(wcp_w P).-2} & The exclusive upper bound on \coq{s}, giving \(w-2\) neighboring hybrid steps. \\
\coq{thf_grind P} & The digest-and-counter hash, packaged for a tweakable-hash game. \\
\coq{coll_grind P} & Access to the chain hash and public-key compression while the digest-and-counter hash is being challenged. \\
\coq{wcp_M P} & The type of the digest signed by the scheme. \\
\coq{grind_enc P} & The encoding of a digest and counter as an input to the digest-and-counter hash. \\
\coq{grind_prp P} & The predicate that accepts a hash output when its base-\(w\) digits have the required sum. \\
\coq{wcp_c_max P} & The exclusive upper bound in the signer's counter search. \\
\coq{1} & A target budget of one for the digest-and-counter and public-key-compression games. \\
\coq{thf_pkcompress P} & The public-key-compression function, packaged for a tweakable-hash game. \\
\coq{coll_pkcompress P} & Access to the chain hash and digest-and-counter hash while public-key compression is being challenged. \\
\coq{tcr_budget (wcp_chain_count P) (wcp_w P)} & An upper bound on the target pairs recorded by the chain-collision reduction; its definition is \(\mathit{chain\_count}+\mathit{chain\_count}\cdot w\). \\
\end{longtable}
}

The budget is \coq{wcp_chain_count P} in the \prf, \smud, and \smpre games; it is one in the \stcrp and public-key-compression games; and it is \coq{tcr_budget (wcp_chain_count P) (wcp_w P)} in the chain-collision game.
In the \prf game, this value bounds distinct oracle inputs.
In the tweakable-hash games, it bounds calls that create challenge targets.
It does not bound all evaluations of the public hash functions; \cref{sec:games} explains the available oracles and their query rules.

The table below separates the six cryptographic tasks on the right-hand side.
The \prf and \smud rows each account for more than one term because the proof uses those games at more than one hybrid transition.

\begin{table}[H]
\centering
\caption{The security task in each assumption-game instance.}
\label{tab:bound}
\small
\begin{tabular}{@{}L{0.22\textwidth}L{0.18\textwidth}L{0.52\textwidth}@{}}
\toprule
Game instance & Reduction & Security task \\
\midrule
\prf for key derivation & \coq{R_prf} & Distinguish chain starts derived using one secret seed from independently sampled chain starts. \\
\smud for the chain hash & \coq{R_ud} & Distinguish the hash of a hidden random chain value from an independent random chain value. \\
\stcrp for the digest-and-counter hash & \coq{R_grind} & Given the hash output selected by the signer's counter search, find a different digest and a counter with the same output. \\
\smtcr for public-key compression & \coq{R_compress} & Given the compression of one endpoint vector, find a different endpoint vector with the same compressed value. \\
\smtcr for the chain hash & \coq{R_chain_tcr} & Given one chain-hash input and output, find a different input with the same output at the same chain position. \\
\smpre for the chain hash & \coq{R_chain_pre} & Given the chain-hash output of a hidden random input, find any input with that output at the same chain position. \\
\bottomrule
\end{tabular}
\end{table}

\section{Review the exact assumption games}
\label{sec:games}

The abstract theorem reduces forgery to the exact games defined in \pathcode{proof/ssprove/lib/games/}.
A cryptographic review must establish whether those games match the assumptions used by an external security analysis.
These expanded names are mnemonic; the oracle definitions determine the assumptions.
For that comparison, each subsection below first states the assumption as a mathematical experiment and explains its cryptographic purpose.
It then compares that experiment with the corresponding \SSProve game.

\subsection{The hash functions supplied to the games}

The games are generic over the hash function being challenged.
The records in \cref{lst:hash-records} supply the function's types, evaluation function, and the distributions from which the game samples keys, inputs, or outputs.

\lstinputlisting[
  caption={The hash-function records used to instantiate every assumption game.},
  label={lst:hash-records},
  firstline=69,lastline=99,firstnumber=69
]{../proof/ssprove/lib/games/thf.v}

The record \coq{khf} describes the keyed function used to derive chain starts.
Its fields specify a key type, an input type, an output type, the evaluation function, and samplers for keys and outputs.

The record \coq{thf} describes a tweakable hash \(H(P,t,x)\).
The value \(P\) is the public parameter, the tweak \(t\) is an encoded address, and \(x\) is the hash input.
The fields \coq{th_dPP}, \coq{th_dM}, and \coq{th_dN} specify how a game samples a public parameter, a hidden input, and an output for the ideal game.
The record itself imposes no security property on \coq{th_eval}.

A reduction that attacks one tweakable hash must still simulate the other hash functions used by the signature scheme.
Those functions share the same public parameter, which remains hidden while the reduction creates its targets.
The record \coq{thf_coll} therefore supplies a \emph{collection oracle} that evaluates the other functions under the game's public parameter.
The security-relevant choices are therefore the challenged evaluation function, the samplers used by its game, and the other functions available through the collection oracle.

\lstinputlisting[
  caption={The three functions supplied to the tweakable-hash games.},
  label={lst:hash-instances},
  firstline=145,lastline=160,firstnumber=145
]{../proof/ssprove/wots/games.v}

All three \coq{th_eval} definitions convert the game's encoded tweak to the structured address expected by the corresponding \WOTSC function using \coq{decode_tweak_core}.
The equalities proved on lines 231 to 240 of the same file show that, for every structured address \coq{a}, evaluation at \coq{encode_tweak P a} is exactly \coq{wcp_F}, \coq{wcp_T_sf}, or \coq{wcp_H_grind} at \coq{a}.
These equalities establish that the generic game records challenge the intended functions at the intended addresses.

\Cref{tab:hash-instances} separates the fields that affect the games from fields present only because every \coq{thf} record has the same shape.

\begin{table}[H]
\centering
\caption{Security-relevant fields of the three tweakable-hash instances.}
\label{tab:hash-instances}
\small
\begin{tabular}{@{}L{0.20\textwidth}L{0.29\textwidth}L{0.43\textwidth}@{}}
\toprule
Instance & Challenged function & Samplers used by its games \\
\midrule
\coq{thf_chain} & The chain hash \(F(P,t,x)\). & \coq{th_dPP} uses \coq{wcc_dPP H}; \coq{th_dM} and \coq{th_dN} independently sample chain values using \coq{wcc_dN H}. \\
\coq{thf_pkcompress} & Public-key compression \(T_{\mathsf{sf}}(P,t,xs)\). & Only \coq{dPP} is used by the instantiated \smtcr game; its input and ideal-output samplers are unused. \\
\coq{thf_grind} & The digest-and-counter hash \(H_{\mathsf{grind}}(P,t,(m,c))\). & Only \coq{dPP} is used by the instantiated \stcrp game; its input and ideal-output samplers are unused. \\
\bottomrule
\end{tabular}
\end{table}

For \coq{thf_chain}, lines 146 to 148 set both the input and output type to a chain value and instantiate \coq{th_eval} with the chain hash.
The assignments \coq{th_dM := wcc_dN H} and \coq{th_dN := wcc_dN H} make the hidden input and ideal output independent samples from the same chain-value distribution.
The \smud game uses both samplers, while the \smpre game uses \coq{th_dM} to sample its hidden preimage.

For public-key compression, lines 151 to 153 use a list of chain values as the input and a single chain value as the output.
For the digest-and-counter hash, lines 156 to 160 use a digest-counter pair as the input and a natural number as the output.
The empty-list sampler in the compression record and the samplers that return \((m,0)\) and zero in the digest-and-counter record are placeholders required by the generic record type.
The compression output sampler is also unused.
None of these samplers affects the stated theorem because the corresponding \smtcr and \stcrp games never call them.

The collection definitions on lines 162 to 184 contain mostly type routing.
Their cryptographic content is the mapping in \cref{tab:collections}.

\begin{table}[H]
\centering
\caption{Functions available through each collection oracle.}
\label{tab:collections}
\small
\begin{tabular}{@{}L{0.25\textwidth}L{0.25\textwidth}L{0.42\textwidth}@{}}
\toprule
Challenged function & Collection & Functions available to the reduction \\
\midrule
Chain hash & \coq{coll_chain} & Public-key compression and the digest-and-counter hash. \\
Public-key compression & \coq{coll_pkcompress} & The chain hash and the digest-and-counter hash. \\
Digest-and-counter hash & \coq{coll_grind} & The chain hash and public-key compression. \\
\bottomrule
\end{tabular}
\end{table}

Every function in a collection uses the same public parameter \(P\) as the challenged function.
No collection exposes the challenged function itself.
The \coq{inl} and \coq{inr} branches in the source only select which of the two companion functions to call.

The \stcrp game also needs the encoding of its digest-counter input and the predicate used by the counter search.

\lstinputlisting[
  caption={The input encoding and acceptance predicate for the digest-and-counter game.},
  label={lst:grind-instance},
  firstline=186,lastline=192,firstnumber=186
]{../proof/ssprove/wots/games.v}

The encoding \coq{grind_enc_core} is simply the pair \((m,c)\).
The predicate \coq{grind_prp_core} converts the hash output to \(\ell\) base-\(w\) digits and accepts exactly when those digits are valid and have the required sum.
These are the same encoding and acceptance condition used by the signer's counter search.

\subsection{The protocol shared by the tweakable-hash games}

Let
\[
  H : \mathcal P\times\mathcal T\times\mathcal M\longrightarrow\mathcal N
\]
be a tweakable hash function with samplers \(D_P\), \(D_M\), and \(D_N\) for public parameters, hidden inputs, and ideal outputs.
In this section, a \emph{target} is one hash challenge associated with a tweak \(t\) and created before the adversary learns \(P\).
It does not refer to a \WOTSC public key or forged digest.
The challenge value is later distinguished, collided with, or inverted, depending on the assumption.
\Cref{tab:targets} states exactly what one target contains in each game.

\begin{table}[H]
\centering
\caption{Meaning of one target in each tweakable-hash game.}
\label{tab:targets}
\small
\begin{tabular}{@{}L{0.16\textwidth}L{0.76\textwidth}@{}}
\toprule
Game & Target created before \(P\) is revealed \\
\midrule
\smud & The adversary supplies a fresh tweak \(t\). The game returns either \(H(P,t,x)\) for a hidden \(x\leftarrow D_M\), or an independent \(y\leftarrow D_N\). \\
\smtcr & The adversary supplies a fresh tweak \(t\) and an input \(m\). The game returns \(H(P,t,m)\) and records \((t,m)\). \\
\stcrp & The adversary supplies a fresh tweak \(t\) and a digest \(m\). The game selects the first acceptable counter \(c\), returns \((H(P,t,\mathsf{enc}(m,c)),c)\), and records \((t,m,c)\). \\
\smpre & The adversary supplies a fresh tweak \(t\). The game samples and records a hidden input \(x\leftarrow D_M\), then returns \(H(P,t,x)\). \\
\bottomrule
\end{tabular}
\end{table}

The qualifier \emph{single-function, multi-target} describes the common experiment used by the four tweakable-hash assumptions.
One public parameter \(P\leftarrow D_P\) is shared by at most \(p\) targets, and every target uses a distinct tweak.

Before the public parameter is revealed, the adversary may create targets through the game's target oracle.
The target oracle rejects calls after the reveal, calls beyond the target budget \(p\), and repeated tweaks.
The reveal oracle then returns the same public parameter used for every target and prevents any further target queries.
This implements a two-stage adversary \(\A=(\A_1,\A_2)\): \(\A_1\) creates the targets without knowing \(P\), and \(\A_2\) receives \(P\) and the state retained by \(\A_1\).
In the three search games, a later check names one of these targets by its index and tests a candidate against the recorded data.
The \smud game instead asks the adversary to distinguish the returned target values from independent samples.

Every game samples its public parameter at most once and stores it for later calls.
Lazy sampling changes when the value is drawn, but all oracles in one execution use the same value.
The collection oracle is available before and after the reveal and also uses this value.
We write \(\mathcal C_P\) for this oracle in the mathematical definitions below.
It evaluates the companion hash functions that share \(P\), but it does not evaluate the challenged function \(H\).

The target, reveal, and phase checks occur before any branch on the game bit.
An invalid call therefore aborts in both games at the same point.
This placement matters under the subprobability semantics from \cref{sec:ssprove-semantics}: an assertion present on only one side could create a distinguishing signal unrelated to the intended collision or preimage event.

The pen-and-paper search experiments let \(\A_2\) output one candidate and define a predicate \(W\) that decides whether it wins.
The \SSProve games expose \(W\) as a \coq{CHECK} oracle: the real game returns \(W\), while the ideal game returns \coq{false}.
A package may call \coq{CHECK} more than once, but after the reveal \(W\) can be computed from \(P\), the selected target's call and reply, and the submitted candidate.
The package can therefore select one candidate for which the same predicate holds.
This correspondence concerns the experiment's outputs only; a concrete cost analysis must also count the hash evaluations used to test candidates.

The target budget \(p\) counts target-oracle calls only.
The collection oracle has no query bound, and after the reveal an adversary can evaluate the public challenged function locally.
Consequently, \(p\) is not a bound on the adversary's total number of hash evaluations.
Any external concrete analysis that uses a total hash-query bound must account separately for these evaluations and for those performed by the reduction.

\subsection{The PRF game}

Let \(F:\mathcal K\times\mathcal M\to\mathcal N\) be a keyed function.
Its standard PRF advantage for an adversary making at most \(q\) queries is
\begin{equation}
\label{eq:prf-definition}
  \Adv^{\mathsf{PRF}}_{F,q}(\A)
  =
  \left|
    \Pr_{k\leftarrow D_K}\!\left[\A^{F(k,\cdot)}=1\right]
    -
    \Pr_{G\leftarrow\mathsf{Func}(\mathcal M,\mathcal N)}\!\left[\A^{G(\cdot)}=1\right]
  \right|,
\end{equation}
where \(G\) is a uniformly sampled function from \(\mathcal M\) to \(\mathcal N\).

\paragraph{Intuition.}
The adversary must distinguish evaluations under one unknown key from the outputs of a random function.
In the \WOTSC reduction, the unknown key is \(\SKseed\), each input identifies one chain, and each output is the secret value at the start of that chain.
The assumption justifies replacing all derived chain starts with independent random values.

\lstinputlisting[
  caption={The complete \prf game used for secret-seed key derivation.},
  label={lst:prf-game},
  firstline=48,lastline=77,firstnumber=48
]{../proof/ssprove/lib/games/prf.v}

\paragraph{Correspondence with the \SSProve game.}
The game exports one oracle, \coq{FQUERY}.
Lines 56 to 62 limit the interaction to \(q\) distinct inputs and record each input before the game branches.
The location \coq{prf_k_loc} in the real branch stores only the secret key; it does not store input-output pairs.
Lines 64 to 72 sample this key on the first query, reuse it for every later query, and return \(F(k,m)\).
The list read on line 56 records the queried inputs solely to enforce the distinct-input condition on line 58.

Lines 73 to 75 sample a fresh output from \coq{kh_dN} for every accepted ideal query.
Every accepted query has a new input, so this is lazy sampling of a random function at previously undefined points.
For two distinct inputs \(m_1\) and \(m_2\), a uniformly random function assigns independent uniform values to \(m_1\) and \(m_2\); those values may still coincide.
When \coq{kh_dN} is the uniform sampler on \(\mathcal N\), the ideal branch has exactly this distribution.
A repeated input is rejected on line 58 before either branch executes, so it cannot reveal the difference between the branches.
Without this assertion, querying the same input twice and comparing the replies would distinguish the two branches exactly as described.

Equation~\eqref{eq:prf-definition} normally permits repeated inputs and requires the random function to answer them consistently.
An adversary for the standard experiment can cache each first reply itself and answer a repeated query from that cache.
This transformation preserves its transcript because both \(F(k,\cdot)\) and a sampled random function return the same value when queried twice at the same input.
The distinct-input restriction in \coq{PRF_game} is therefore equivalent for this purpose.
The external bound must nevertheless use the same key, input, and output distributions and must cover the number of distinct inputs shown in the theorem.

\subsection{The SM-UD game}

For \(b\in\{0,1\}\), define the target oracle
\[
  \mathcal O^{\mathsf{UD}}_{P,b}(t)=
  \begin{cases}
    H(P,t,x), & b=0,\quad x\leftarrow D_M,\\
    y,        & b=1,\quad y\leftarrow D_N.
  \end{cases}
\]
An \smud adversary creates at most \(p\) targets under distinct tweaks while \(P\) is hidden, receives \(P\), and guesses which oracle it received.
Its advantage is
\begin{equation}
\label{eq:smud-definition}
  \Adv^{\mathsf{SMUD}}_{H,p}(\A)
  =
  \left|
    \Pr\!\left[\A^{\mathcal O^{\mathsf{UD}}_{P,0},\mathcal C_P}=1\right]
    -
    \Pr\!\left[\A^{\mathcal O^{\mathsf{UD}}_{P,1},\mathcal C_P}=1\right]
  \right|.
\end{equation}
Each probability samples \(P\leftarrow D_P\) and includes the two-stage protocol described above.

\paragraph{Intuition.}
The hash of a hidden random input should look like a fresh random output, even after the public parameter is revealed.
In the \WOTSC proof, this property lets a hybrid replace unrevealed chain values with random values without giving the forger a detectable signal.
Distinct tweaks keep the targets in separate hash domains.

\lstinputlisting[
  caption={The complete single-function multi-target undetectability game.},
  label={lst:smud-game},
  firstline=54,lastline=113,firstnumber=54
]{../proof/ssprove/lib/games/smud.v}

\paragraph{Correspondence with the \SSProve game.}
The oracle \coq{UQUERY} accepts a tweak before the public parameter is revealed.
Lines 64 to 74 enforce the phase, target budget, and distinct-tweak conditions.
In the real game, lines 76 to 85 sample the public parameter once, sample a hidden input, and return its hash under the chosen tweak.
In the ideal game, lines 86 to 88 return an independent output from \coq{th_dN}.
The oracle \coq{UGETPP} reveals the public parameter, and \coq{UCOLLQ} evaluates the companion functions under that same parameter.

These clauses implement \eqref{eq:smud-definition} directly.
An external bound must use the same hidden-input and ideal-output distributions and must include the collection access shown here.

\subsection{The SM-TCR game}

For each target query \((t_i,m_i)\), the challenger records the pair and returns
\[
  y_i=H(P,t_i,m_i).
\]
The adversary makes at most \(p\) such queries under distinct tweaks before learning \(P\).
After receiving \(P\), it outputs \((j,m^\star)\) and wins if
\begin{equation}
\label{eq:smtcr-definition}
  m^\star\ne m_j
  \quad\land\quad
  H(P,t_j,m^\star)=H(P,t_j,m_j).
\end{equation}
The adversary has access to \(\mathcal C_P\) throughout the experiment.

\paragraph{Intuition.}
The adversary commits to each target input before it learns the public parameter that selects the hash function.
It must then find a different input with the same output under the target's original tweak.
This is the target-collision, or universal one-way, analogue of collision resistance; the first input is fixed before the hash function becomes known.

\lstinputlisting[
  caption={The complete single-function multi-target target-collision game.},
  label={lst:smtcr-game},
  firstline=51,lastline=124,firstnumber=51
]{../proof/ssprove/lib/games/smtcr.v}

\paragraph{Correspondence with the \SSProve game.}
Before the reveal, \coq{TQUERY} accepts a tweak and an input, records both, and returns their hash output.
After the reveal, \coq{CHECK} accepts a target index and a replacement input.
Lines 96 to 108 return \coq{true} exactly when the index exists, the replacement input differs from the recorded input, and both inputs hash to the same output under the recorded tweak.

The theorem uses this game twice.
In the public-key-compression instance, an input is a vector of chain endpoints, and a successful candidate is a different vector with the same compressed value.
In the chain-hash instance, an input is one chain value, and a successful candidate is a different value with the same hash output at a recorded chain tweak.
The theorem statement supplies a separate target budget for each instance.

The oracle order and winning condition are exactly those of \eqref{eq:smtcr-definition}.
An external target-collision bound must also include the appropriate companion collection for each instance.

\subsection{The S-TCR-P game}

The \stcrp game adds a counter search and an acceptance predicate to the target-collision experiment above.
Let \(\mathsf{enc}:\mathcal M_0\times\mathbb N\to\mathcal M\) encode a digest and counter, and let \(\mathsf{Prop}:\mathcal N\to\{0,1\}\).
On a target query \((t_i,m_i)\), the challenger computes
\[
  c_i=\min\left\{c\in\mathbb N:\mathsf{Prop}\bigl(H(P,t_i,\mathsf{enc}(m_i,c))\bigr)=1\right\}.
\]
The definition presumes that this set is nonempty.
The challenger records \((t_i,m_i,c_i)\) and returns
\[
  \left(H(P,t_i,\mathsf{enc}(m_i,c_i)),c_i\right).
\]
After \(P\) is revealed, the adversary outputs \((j,m^\star,c^\star)\) and wins if
\begin{equation}
\label{eq:stcrp-definition}
  m^\star\ne m_j
  \quad\land\quad
  H(P,t_j,\mathsf{enc}(m_j,c_j))
  =H(P,t_j,\mathsf{enc}(m^\star,c^\star)).
\end{equation}
The target queries use distinct tweaks and are limited to \(p\).

\paragraph{Intuition.}
This is target-collision resistance for a hash input selected by the signer's counter search.
The target digest is fixed before \(P\) is known, but its counter depends on \(P\) because the challenger chooses the first acceptable hash output.
The adversary must replace the digest while preserving the recorded output under the same tweak.

\lstinputlisting[
  caption={The complete predicate-constrained target-collision game.},
  label={lst:stcrp-game},
  firstline=71,lastline=158,firstnumber=71
]{../proof/ssprove/lib/games/stcrp.v}

\paragraph{Correspondence with the \SSProve game.}
The oracle \coq{SQUERY} searches counters in increasing order, starting at zero and stopping before \(\mathsf{cap}\).
It selects the first counter whose hash output satisfies \coq{prp}.
If the search succeeds, the oracle records the tweak, base digest, and selected counter and returns the hash output and counter.
If no counter succeeds, lines 102 to 110 abort.

After \coq{SGETPP} reveals the public parameter, \coq{SCHECK} accepts an index, a replacement digest, and a replacement counter.
Lines 129 to 142 return \coq{true} exactly when the index names a recorded target, the replacement digest differs from the target digest, and the two encoded digest-counter pairs have the same hash output under the target tweak.
The replacement counter has type \coq{nat} and is not required to be below \(\mathsf{cap}\).
Allowing any natural-number replacement counter gives the adversary at least as many choices as a game that restricts the counter to the signer's range.

The \WOTSC reduction uses one target and instantiates \coq{enc}, \coq{prp}, and \coq{cap} with the signer's encoding, acceptance predicate, and counter bound.
The pen-and-paper definition searches all natural-number counters, whereas \coq{SQUERY} searches only below \(\mathsf{cap}\) and aborts if that finite search fails.
On executions where the bounded search succeeds, both experiments select the same counter.
The abort removes the remaining executions from both branches of the \SSProve game, so it cannot increase the distinguishing advantage.
An external bound for the unbounded experiment can therefore cover the bounded experiment if it also covers the unrestricted replacement counter and the collection oracle.

\subsection{The SM-PRE game}

For each target query on a tweak \(t_i\), the challenger samples \(x_i\leftarrow D_M\), keeps \(x_i\) hidden, and returns
\[
  y_i=H(P,t_i,x_i).
\]
The adversary makes at most \(p\) such queries under distinct tweaks before learning \(P\).
After receiving \(P\), it outputs \((j,m^\star)\) and wins if
\begin{equation}
\label{eq:smpre-definition}
  H(P,t_j,m^\star)=y_j.
\end{equation}
The adversary has access to \(\mathcal C_P\) throughout the experiment.

\paragraph{Intuition.}
The adversary sees the hash of a random hidden chain value and must find any preimage after the public parameter is revealed.
Unlike second-preimage resistance, this experiment does not require \(m^\star\ne x_j\) because the target input was never given to the adversary.
In the \WOTSC proof, this property covers a forgery that moves backward through a hash chain to a hidden value.

\lstinputlisting[
  caption={The complete single-function multi-target preimage game.},
  label={lst:smpre-game},
  firstline=53,lastline=127,firstnumber=53
]{../proof/ssprove/lib/games/smpre.v}

\paragraph{Correspondence with the \SSProve game.}
Before the reveal, \coq{PQUERY} accepts a tweak, samples an input from \coq{th_dM}, records the input without revealing it, and returns its hash output.
After the reveal, \coq{PCHECK} accepts a target index and a candidate input.
Lines 99 to 111 return \coq{true} exactly when the candidate hashes to the recorded output under the recorded tweak.
The game does not require the candidate to differ from the hidden input because the adversary was never given that input.

These clauses implement \eqref{eq:smpre-definition} directly.
An external preimage bound must use the same hidden-input distribution and include the collection access.

\subsection{What the published generic bounds provide}
\label{sec:generic-bounds}

H\"ulsing and Kudinov collect the generic bounds in the quantum random-oracle model (QROM) relevant to the games above~\cite[Table~1]{HulsingKudinov2022}.
Their parameter \(q\) counts quantum queries to the modeled function, \(p\) counts classical target-oracle queries, and \(n\) is the output length.

\begin{table}[H]
\centering
\small
\caption{Published QROM bounds relevant to the formal games.}
\label{tab:published-bounds}
\begin{tabular}{L{2.0cm}L{6.1cm}L{6.5cm}}
\toprule
Formal game & Published bound & Status \\
\midrule
\prf
  & \(\Adv^{\mathsf{PRF}}\leq 12q/\sqrt{2^n}\)
  & Proved for a random keyed-function family by H\"ulsing and Kudinov, Lemma~2~\cite{HulsingKudinov2022}. \\
\smud
  & \(\Adv^{\mathsf{SM\text{-}UD}}\leq 12q/\sqrt{2^n}\)
  & Proved for the standalone game by H\"ulsing and Kudinov, Lemma~4~\cite{HulsingKudinov2022}. \\
\smtcr
  & \(\operatorname{Succ}^{\mathsf{SM\text{-}TCR}}\leq 8(2q+1)^2/2^n\)
  & Proved for the simple quantum-random-oracle construction by Bernstein et al., Theorem~11~\cite{BernsteinEtAl2019Sphincs}; the underlying generic second-preimage analysis appears in H\"ulsing, Rijneveld, and Song~\cite{HulsingRijneveldSong2016}. \\
\smpre
  & \(\operatorname{Succ}^{\mathsf{SM\text{-}PRE}}\leq O((q+1)^2/2^n)\)
  & The tight multi-target bound for a length-preserving function depends on a conjectured multi-target decisional second-preimage bound; an unconditional reduction loses a factor \(p\)~\cite{BernsteinHulsing2019DSPR,HulsingKudinov2022}. \\
\stcrp
  & \(\operatorname{Succ}^{\mathsf{S\text{-}TCR\text{-}P}}\leq \dfrac{32(q+1/2)^2}{2^n}+\dfrac{16q^2}{|\mathcal P|}\)
  & The paper calls the fixed-digit-sum instance \(\mathsf{S\text{-}TCR}(+C)\) and proves the displayed bound for that instance~\cite{hulsing2023sphincs+}. \\
\bottomrule
\end{tabular}
\end{table}

The cited papers use advantage for the two distinguishing experiments and success probability for the three collision or preimage experiments.
The notation in the table follows that distinction.

The cited papers prefix the \smud, \smtcr, and \smpre names with \(\mathsf{W}\) to denote the target-before-reveal order already specified above.

The table does not provide values for the six formal \(\varepsilon\)-terms because the published and formal models differ in several respects.
The published bounds are oracle-model statements.
In particular, the simple-construction \smtcr theorem does not give the first-stage adversary direct access to the underlying random oracle~\cite[Theorem~11]{BernsteinEtAl2019Sphincs}.
In the concrete formal instances, the hash maps are ordinary Rocq functions that an adversary package may use in its local computations on known inputs.
The formal model does not count these local evaluations.
An instantiation therefore needs an oracle simulation and a cost model that accounts for both exported oracle calls and local hash evaluations.

The formal \smud, \smtcr, and \smpre games also include collection oracles.
Definition~7 of H\"ulsing and Kudinov gives the first-stage adversary a collection oracle under the same hidden public parameter, but it requires the challenge tweaks and collection-query tweaks to be disjoint~\cite{HulsingKudinov2022}.
The formal games do not enforce this disjointness, permit collection queries in both stages, and do not bound their number.
H\"ulsing and Kudinov conjecture that their standalone generic bounds remain unchanged for collections; they do not prove that extension.

The formal chain hash has another mismatch: it maps an \(n\)-bit input to an \(n\)-bit output.
The direct multi-target preimage theorem of H\"ulsing, Rijneveld, and Song instead assumes a compressing function whose input length is \(cn\) for a constant \(c>1\), together with \(p=o(2^n)\)~\cite{HulsingRijneveldSong2016}.
For length-preserving functions, H\"ulsing and Kudinov derive \smpre from \smtcr and single-function multi-target decisional second-preimage resistance (\(\mathsf{SM\text{-}DSPR}\)).
Bernstein and H\"ulsing bound the one-target random-function DSPR advantage by \(32q^2/2^n\) and give a reduction to multiple targets that loses a factor equal to the number of targets~\cite[Theorems~30 and~33]{BernsteinHulsing2019DSPR}.
The target-independent multi-target DSPR bound used to obtain the tight \(O((q+1)^2/2^n)\) preimage bound remains a conjecture~\cite{BernsteinEtAl2019Sphincs,HulsingKudinov2022}.

The published \stcrp bound does not apply directly to the formal game.
This mismatch does not give a known attack against the formal game.
It means that a separate argument is needed before the published bound can be used for \(\epsgr\).

The formal target oracle selects the least acceptable counter below a finite cap and aborts if no such counter exists.
The published experiment assumes that an acceptable natural-number counter exists, and its generic analysis uses an unbounded search.
The bounded search matches the function analyzed by the formal proof, so replacing it with the paper's unbounded search would make the formal game less faithful.
A correspondence proof should instead relate collision success in the bounded game to collision success in the published experiment.

The formal interface also exposes \coq{SCOLLQ}, which evaluates the chain hash or public-key-compression hash under the hidden public parameter.
The published \stcrp experiment has no collection oracle.
This extra oracle makes the formal assumption stronger, but the grinding reduction \coq{R_grind} never calls it.
Removing \coq{SCOLLQ} from the \stcrp interface would therefore weaken the assumption without changing that reduction.

Finally, the formal interface permits repeated \coq{SCHECK} calls, whereas the paper's second-stage adversary returns one candidate.
After the public parameter is known, the adversary can evaluate the collision predicate itself.
A correspondence proof can use one successful candidate, but its QROM query count must include the evaluations used to find that candidate.

Using the published bound for \(\epsgr\) therefore requires this game-correspondence argument and a model relating the formal hash functions to the random tweakable hash analyzed in the paper.

The table identifies the relevant external results but does not justify substituting their bounds for the six \(\varepsilon\)-terms.
Its \(q\) counts function or random-oracle queries.
The formal games bound only the number of target queries and place no cost semantics on the adversary or the reduction.
The cited quantum bounds also apply to classical adversaries that obey the same query limit, but they do not add quantum or QROM semantics to the \SSProve development.
Instantiating the formal axioms requires the missing experiment correspondences and resource analysis.

\section{Conclusion}

The Rocq kernel checks a classical reduction from any \coq{ValidPackage} adversary in the known-message \WOTSC experiment to the six assumption-game adversaries constructed by the proof.
It also checks the advantage inequality in \cref{lst:abstract-theorem}.

The checked reduction is not yet a concrete QROM security theorem.
The formal proof has no resource semantics, and no checked argument connects all six formal games to the cited QROM experiments.
The tight published bound relevant to the length-preserving \smpre game is conjectural, and the formal \stcrp game differs from the published experiment as described in \cref{sec:generic-bounds}.

\bibliographystyle{alpha}
\bibliography{SSPROVE_SECURITY_REVIEW}

\end{document}
